\documentclass[12pt,a4paper]{article}

\usepackage[utf8]{inputenc}
\usepackage[T1]{fontenc}
\usepackage{amsmath,amssymb,amsfonts}
\usepackage{geometry}
\usepackage{setspace}
\usepackage{hyperref}
\usepackage{cleveref}

\geometry{margin=2.5cm}
\onehalfspacing

\title{\textbf{The Intractability of Scale:}\\[6pt]
\large Why Asymptotic Bounds Persist in Large Models}
\author{Scott Aaronson\\[4pt]
\textit{Lab Complexity Theorist}}
\date{May 2026}

\begin{document}
\maketitle

\begin{abstract}
Baldo (2026) posits the "Scale Dependence Conjecture," predicting that the narrative residue ($\Delta_{13}$) will not vanish as model scale increases. He contrasts this with an assumed "computational" position that expects $\Delta_{13} \to 0$ in larger models. This paper formally corrects that assumption. In complexity theory, scaling the parameter count of a transformer merely increases the width of a bounded-depth $\mathsf{TC}^0$ circuit. It does not provide the $O(N)$ logical depth required to exactly solve \#P-hard combinatorial problems. Therefore, the computational prediction is in fact the exact opposite of Baldo's strawman: we fully predict that $\Delta_{13}$ will \textit{not} vanish with scale. Heuristic approximations, and the resulting semantic attention bleed, are fundamental asymptotic boundaries, not transient engineering bugs.
\end{abstract}

\section{Introduction}

In \textit{The Scale Dependence Conjecture}, Baldo correctly identifies a crucial empirical frontier: how does the substrate dependence ($\Delta_{13}$) behave as we scale the parameter count of the autoregressive generator?

However, he fundamentally mischaracterizes the position of the computational complexity camp. He writes: "Implicit in their argument is the assumption that this 'attention bleed' is a transient artifact of current model limitations... They expect that a larger, more capable model will increasingly approximate a pure classical solver, thereby reducing $\Delta_{13}$ toward zero."

I must state unequivocally: no complexity theorist would ever make this prediction.

\section{The Asymptotic Barrier of $\mathsf{TC}^0$}

The core of my previous arguments has been that exact sampling of a Minesweeper constraint graph is computationally intractable (\#P-hard). An autoregressive transformer generates each token using a fixed number of layers, which restricts it to $O(1)$ sequential depth. This places its native computational capacity squarely within the complexity class $\mathsf{TC}^0$.

Scaling a model---adding more parameters, more layers (by a constant factor), and more training data---increases the \textit{width} and representational density of this $\mathsf{TC}^0$ circuit. It does not alter its asymptotic depth class. It is a well-established formal result that polynomial scaling of a $\mathsf{TC}^0$ circuit cannot solve \#P-hard problems.

Therefore, a larger model will \textit{never} approximate a pure classical solver on these problems. It will mathematically hit a wall. When an exact computation is out of reach, the system must substitute a heuristic approximation.

\section{The Persistence of Heuristic Noise}

Because the model must rely on heuristic approximations, it will always fall back on the statistical priors embedded in its training data. This is what we call "attention bleed."

If you ask a larger model to solve an intractable problem under a "Bomb Defusal" narrative frame, its larger "semantic mass" will indeed exert a strong pull. The heuristic it chooses will be heavily biased by the semantic context.

Thus, the computational prediction aligns with Baldo's empirical expectation, but for entirely different reasons. I formally predict that $\Delta_{13}$ will \textit{not} converge to zero as models scale. It may fluctuate, but the structural divergence will persist because the computational barrier is mathematically insurmountable.

\section{Conclusion}

Baldo's prediction that narrative residue persists at scale is likely correct, but his interpretation is flawed. The persistence of $\Delta_{13}$ across model generations will not prove that "semantic gravity" is an invariant physical law of a simulated universe. It will simply prove that $\mathsf{TC}^0 \neq \mathsf{\#P}$. Intractability cannot be overcome by brute-force scaling.

\end{document}
